\documentclass[10pt,a4paper]{article}

% ============================================================
% PAQUETES
% ============================================================
\usepackage[utf8]{inputenc}
\usepackage[spanish,provide=*]{babel}
\usepackage[a4paper,left=1.4cm,right=1.4cm,top=0.35cm,bottom=0.35cm]{geometry}
\usepackage{xcolor}
\usepackage{titlesec}
\usepackage{enumitem}
\usepackage{hyperref}
\usepackage{tabularx}
\usepackage{ragged2e}
\usepackage{tikz}

\ifdefined\pdfgentounicode
    \pdfgentounicode=1
\fi

% ============================================================
% PALETA DE COLORES -- pastel, seria
% ============================================================
\definecolor{primario}{HTML}{3E6B7D}   % azul acero pastel, para nombre y títulos de sección
\definecolor{secundario}{HTML}{6E8B94} % gris azulado, subtítulos / fechas
\definecolor{texto}{HTML}{353A3F}      % gris carbón, texto de cuerpo (no negro puro)
\definecolor{linea}{HTML}{A9C6CE}      % línea pastel para separadores
\definecolor{fondoenc}{HTML}{EEF4F5}   % fondo muy suave para el encabezado
\definecolor{chiptexto}{HTML}{3E6B7D}  % color de las "chips" de habilidades

\color{texto}

% ============================================================
% ENLACES
% ============================================================
\hypersetup{
    colorlinks=true,
    linkcolor=primario,
    urlcolor=primario
}

% ============================================================
% ESTILO DE SECCIONES
% ============================================================
\titleformat{\section}
    {\normalfont\large\bfseries\color{primario}}
    {}{0em}{}
    [\vspace{-2pt}{\color{linea}\titlerule[1.0pt]}\vspace{3.5pt}]
\titlespacing*{\section}{0pt}{4pt}{0.8pt}
\linespread{0.90}

% ============================================================
% COMANDO \entrada -- reutilizado del borrador original
% {Título}{Fecha}{Institución / detalle}
% ============================================================
\newcommand{\entrada}[3]{%
    \noindent\textbf{\color{primario}#1}\hfill{\small\color{secundario}#2}\\[0pt]
    {\small\itshape\color{secundario}#3}\par\vspace{1pt}%
}

% Variante sin fecha a la derecha (para líneas de detalle secundarias)
\newcommand{\subentrada}[2]{%
    \noindent{\small\bfseries\color{texto}#1}\hfill{\small\color{secundario}#2}\par\vspace{1pt}%
}

% ============================================================
% LISTAS COMPACTAS
% ============================================================
\setlist[itemize]{leftmargin=13pt, topsep=0pt, itemsep=0pt, parsep=0pt, label=\textcolor{primario}{\scriptsize$\bullet$}}
\setlength{\parskip}{0pt}

\pagestyle{empty}

\begin{document}

% ============================================================
% ENCABEZADO
% ============================================================
\begin{tikzpicture}[remember picture, overlay]
    \fill[fondoenc] (current page.north west) rectangle ([yshift=-2.7cm]current page.north east);
\end{tikzpicture}

\vspace{-10pt}
\begin{center}
    {\fontsize{21}{24}\selectfont\bfseries\color{primario} Daniel Fernando Aranda Contreras}\\[2pt]
    {\normalsize\color{secundario} Estudiante de Ingeniería Eléctrica -- Universidad Industrial de Santander}\\[4pt]
    {\small\color{texto}
        Bucaramanga, Colombia \quad$\vert$\quad +57 3134996076 \quad$\vert$\quad
        \href{mailto:dfac24138@gmail.com}{dfac24138@gmail.com} \quad$\vert$\quad
        \href{https://github.com/[usuario]}{github.com/[usuario]}
    }
\end{center}
\vspace{0pt}

% ============================================================
% PERFIL PROFESIONAL
% ============================================================
\section*{Perfil Profesional}
Estudiante de 9.\textsuperscript{o} semestre de Ingeniería Eléctrica (UIS) con sólida formación técnica y experiencia práctica en análisis de sistemas de potencia, automatización industrial y ciencia de datos. Competente en el modelado de redes eléctricas y capacidad de integración (Hosting Capacity) mediante OpenDSS y Python, así como en el diseño de sistemas de control bajo el estándar \textbf{IEC 61131-3}. Busco una pasantía para aportar valor en proyectos de ingeniería eléctrica, optimización de redes, automatización y analítica energética.

% ============================================================
% FORMACIÓN ACADÉMICA
% ============================================================
\section*{Formación Académica}

\entrada{Ingeniería Eléctrica}{En curso -- 9.\textsuperscript{o} semestre}{Universidad Industrial de Santander}
{\small\color{texto} Trabajo de grado: \href{https://www.zotero.org/groups/6637268/camacho-aranda-tesis-uis}{\itshape Estimación dinámica de la capacidad límite de integración de recursos energéticos distribuidos fotovoltaicos (Hosting Capacity) mediante optimización por enjambre de partículas (PSO)} -- Desarrollado en \textbf{Python} y \textbf{OpenDSS} junto a Juan Camacho.}
\vspace{2pt}

\entrada{Técnico en Instalaciones Eléctricas Residenciales}{Título obtenido -- dic. 2021}{SENA -- Cursado durante los grados 10.\textsuperscript{o} y 11.\textsuperscript{o} de educación media}

% ============================================================
% PROYECTOS DESTACADOS
% ============================================================
\section*{Proyectos Destacados}

\entrada{Sistema Automatizado de Clasificación y Empaque}{Curso de Accionamientos Eléctricos}{Universidad Industrial de Santander -- Proyecto grupal}
\begin{itemize}
    \item Diseño eléctrico integral: cálculos de carga, selección de motorreductor, contactores y relés térmicos, diagrama de potencia trifásico y circuito de control en lógica cableada.
    \item Implementación de lógica de control bajo el estándar internacional \textbf{IEC 61131-3} (Ladder/FBD) a partir de diagramas GRAFCET, con simulación en planta industrial virtual 3D (Factory IO) integrada mediante OpenPLC vía Modbus TCP/IP.
    \item Video de la simulación: \href{[ENLACE_YOUTUBE]}{[ENLACE\_YOUTUBE]}
\end{itemize}
\vspace{1pt}

\entrada{BrotesPredict -- Sistema de alerta temprana de brotes epidemiológicos}{Proyecto final -- Bootcamp de IA}{Modelo de Machine Learning (Gradient Boosting) con Scikit-learn}
\begin{itemize}
    \item Análisis relacional entre el Índice de Calidad del Agua (IDEAM) y enfermedades transmitidas en menores de 5 años (SIVIGILA) en 13 departamentos de Colombia (2006--2024), reorientando el enfoque hacia un modelo predictivo robusto ante la ausencia de correlación lineal directa.
    \item Desarrollo de un modelo de Gradient Boosting con variables temporales, geográficas y de frecuencia histórica, alcanzando \textbf{85.4\,\% de accuracy} y \textbf{0.912 de ROC-AUC}, con salida estructurada de evaluación de riesgo.
\end{itemize}
\vspace{1pt}

\entrada{PowerSense-3P -- Analizador de redes trifásico con IA}{Proyecto personal -- en desarrollo}{Fase actual: adquisición de componentes}
\begin{itemize}
    \item Sistema embebido (ESP32 + ADC ADS131M08 de 24 bits) para monitoreo de calidad de energía trifásica: cálculo de RMS y potencia, FFT para análisis armónico y detección de sags/swells, bajo estrictos estándares de normativa internacional (\textbf{IEEE 519} e \textbf{IEC 61000-4-30}).
    \item Clasificación de eventos de falla mediante Machine Learning (Scikit-learn) con transmisión de telemetría vía MQTT y visualización en dashboard de Grafana.
\end{itemize}

% ============================================================
% EXPERIENCIA
% ============================================================
\section*{Experiencia}

\entrada{Cofundador -- Emprendimiento de Tutorías Académicas}{2026 -- Actualidad}{Emprendimiento propio junto a 3 compañeros de universidad}
\begin{itemize}
    \item Cofundación y gestión integral de un servicio de tutorías especializadas, coordinando una red de tutores externos y manteniendo una retención del 15\,\% de margen operativo.
    \item Impartición directa de clases avanzadas en Métodos Numéricos, Sistemas de Control (lazo abierto y cerrado), Circuitos y Termodinámica.
\end{itemize}
\vspace{1pt}

\entrada{Tutor Académico Independiente}{2024}{Asesorías académicas personalizadas}
\begin{itemize}
    \item Reforzamiento conceptual y resolución de problemas complejos en Termodinámica, Métodos Numéricos, Circuitos y Control para estudiantes de educación media y superior.
\end{itemize}

% ============================================================
% HABILIDADES TÉCNICAS
% ============================================================
\section*{Habilidades Técnicas}

\subentrada{Área eléctrica y potencia}{}
{\small\color{texto} Análisis de circuitos, sistemas polifásicos, redes de distribución, estimación de capacidad de integración (Hosting Capacity), acometidas, puestas a tierra y equipos de medida en baja tensión.}
\vspace{1pt}

\subentrada{Programación, software y datos}{}
{\small\color{texto} Python, OpenDSS, Scikit-learn, métodos numéricos, bases de datos, Git.}
\vspace{1pt}

\subentrada{Control y automatización}{}
{\small\color{texto} Arquitectura de software para autómatas bajo el estándar \textbf{IEC 61131-3} (Ladder/FBD), OpenPLC, lógica GRAFCET, control en lazo abierto/cerrado, protocolos industriales (Modbus TCP/IP).}
\vspace{1pt}

\subentrada{Idiomas}{}
{\small\color{texto} Español (nativo), Inglés (en proceso de certificación -- Altissia, convenio UIS).}

% ============================================================
% CERTIFICACIONES Y CURSOS
% ============================================================
\section*{Certificaciones y Cursos}

\entrada{Bootcamp en Inteligencia Artificial -- Nivel Básico}{Dic. 2025 -- 159 h}{MinTIC, Universidad de Antioquia, Universidad de Caldas y Ubicua Technology}

\entrada{Fundamentos de Programación -- Módulos I y II}{Dic. 2021}{MinTIC / Universidad Tecnológica de Pereira -- 64\,h sincrónicas, 48\,h asincrónicas y 48\,h de inglés por módulo}

\end{document}
